\documentclass[12pt,t,handout]{beamer}
\usepackage{graphicx}
\beamerdefaultoverlayspecification{<+->}
\setbeamercovered{transparent}
\setbeameroption{show notes}
\setbeamertemplate{note page}[plain]
\usepackage{listings}
\usepackage{datetime}
\usepackage{url}

%make math easier
\usepackage{bm}
\usepackage{amsmath}
\usepackage{mathtools}
\newcommand{\D}{\mathcal{D}}
\newcommand{\E}{\mathbb{E}}
\newcommand{\F}{\mathcal{F}}
\newcommand{\X}{\mathcal{X}}
\newcommand{\lik}{\mathcal{L}}
\DeclarePairedDelimiterX{\infdivx}[2]{(}{)}{%
  #1\;\delimsize\|\;#2%
}
\newcommand{\infdiv}{D\infdivx}
\DeclarePairedDelimiter{\norm}{\lVert}{\rVert}
\DeclareMathOperator*{\argmin}{arg\,min}
\DeclareMathOperator*{\argmax}{arg\,max}

%Bibliography
\usepackage{natbib}
\bibpunct{(}{)}{,}{a}{}{;}
\usepackage{bibentry}
%\nobibliography*

\def\notescolors{1}
% header.tex: boring LaTeX/Beamer details + macros

% get rid of junk
\usetheme{default}
\beamertemplatenavigationsymbolsempty
\hypersetup{pdfpagemode=UseNone} % don't show bookmarks on initial view


% font
\usepackage{fontspec}
\setsansfont
  [ ExternalLocation = fonts/ ,
    UprightFont = *-regular ,
    BoldFont = *-bold ,
    ItalicFont = *-italic ,
    BoldItalicFont = *-bolditalic ]{texgyreheros}
\setbeamerfont{note page}{family*=pplx,size=\footnotesize} % Palatino for notes
% "TeX Gyre Heros can be used as a replacement for Helvetica"
% I've placed them in fonts/; alternatively you can install them
% permanently on your system as follows:
%     Download http://www.gust.org.pl/projects/e-foundry/tex-gyre/heros/qhv2.004otf.zip
%     In Unix, unzip it into ~/.fonts
%     In Mac, unzip it, double-click the .otf files, and install using "FontBook"

% named colors
\definecolor{offwhite}{RGB}{255,250,240}
\definecolor{gray}{RGB}{155,155,155}

\ifx\notescolors\undefined % slides
  \definecolor{foreground}{RGB}{255,255,255}
  \definecolor{background}{RGB}{24,24,24}
  \definecolor{title}{RGB}{107,174,214}
  \definecolor{subtitle}{RGB}{102,255,204}
  \definecolor{hilit}{RGB}{102,255,204}
  \definecolor{vhilit}{RGB}{255,111,207}
  \definecolor{lolit}{RGB}{155,155,155}
\else % notes
  \definecolor{background}{RGB}{255,255,255}
  \definecolor{foreground}{RGB}{24,24,24}
  \definecolor{title}{RGB}{27,94,134}
  \definecolor{subtitle}{RGB}{22,175,124}
  \definecolor{hilit}{RGB}{122,0,128}
  \definecolor{vhilit}{RGB}{255,0,128}
  \definecolor{lolit}{RGB}{95,95,95}
\fi
\definecolor{nhilit}{RGB}{128,0,128}  % hilit color in notes
\definecolor{nvhilit}{RGB}{255,0,128} % vhilit for notes

\newcommand{\hilit}{\color{hilit}}
\newcommand{\vhilit}{\color{vhilit}}
\newcommand{\nhilit}{\color{nhilit}}
\newcommand{\nvhilit}{\color{nvhilit}}
\newcommand{\lolit}{\color{lolit}}

% use those colors
\setbeamercolor{titlelike}{fg=title}
\setbeamercolor{subtitle}{fg=subtitle}
\setbeamercolor{institute}{fg=lolit}
\setbeamercolor{normal text}{fg=foreground,bg=background}
\setbeamercolor{item}{fg=foreground} % color of bullets
\setbeamercolor{subitem}{fg=lolit}
\setbeamercolor{itemize/enumerate subbody}{fg=lolit}
\setbeamertemplate{itemize subitem}{{\textendash}}
\setbeamerfont{itemize/enumerate subbody}{size=\footnotesize}
\setbeamerfont{itemize/enumerate subitem}{size=\footnotesize}

% page number
\setbeamertemplate{footline}{%
    \raisebox{5pt}{\makebox[\paperwidth]{\hfill\makebox[20pt]{\lolit
          \scriptsize\insertframenumber}}}\hspace*{5pt}}

% add a bit of space at the top of the notes page
\addtobeamertemplate{note page}{\setlength{\parskip}{12pt}}

% default link color
\hypersetup{colorlinks, urlcolor={hilit}}

\ifx\notescolors\undefined % slides
  % set up listing environment
  \lstset{language=bash,
          basicstyle=\ttfamily\scriptsize,
          frame=single,
          commentstyle=,
          backgroundcolor=\color{darkgray},
          showspaces=false,
          showstringspaces=false
          }
\else % notes
  \lstset{language=bash,
          basicstyle=\ttfamily\scriptsize,
          frame=single,
          commentstyle=,
          backgroundcolor=\color{offwhite},
          showspaces=false,
          showstringspaces=false
          }
\fi

% a few macros
\newcommand{\bi}{\begin{itemize}}
\newcommand{\bbi}{\vspace{24pt} \begin{itemize} \itemsep8pt}
\newcommand{\ei}{\end{itemize}}
\newcommand{\ig}{\includegraphics}
\newcommand{\subt}[1]{{\footnotesize \color{subtitle} {#1}}}
\newcommand{\ttsm}{\tt \small}
\newcommand{\ttfn}{\tt \footnotesize}
\newcommand{\figh}[2]{\centerline{\includegraphics[height=#2\textheight]{#1}}}
\newcommand{\figw}[2]{\centerline{\includegraphics[width=#2\textwidth]{#1}}}


%%%%%%%%%%%%%%%%%%%%%%%%%%%%%%%%%%%%%%%%%%%%%%%%%%%%%%%%%%%%%%%%%%%%%%
% end of header
%%%%%%%%%%%%%%%%%%%%%%%%%%%%%%%%%%%%%%%%%%%%%%%%%%%%%%%%%%%%%%%%%%%%%%

% title info
\title{Fair New World}
\subtitle{\scriptsize for the seminar: \textit{Fairness in Machine Learning}, \\
                      organized by M.~Hardt, Fall 2017, UC Berkeley}
\author{\href{http://nimahejazi.org}{Nima Hejazi}
       \\[-10pt]
       }
\institute{Division of Biostatistics \\
           University of California, Berkeley \\
           \href{https://www.stat.berkeley.edu/~nhejazi}
             {\tt \scriptsize \color{foreground}
               stat.berkeley.edu/\textasciitilde{}nhejazi
             }
           \\[4pt]
           \includegraphics[height=20mm]{Figs/seal-berkeley.png}
           \\[-12pt]
          }
\date{
  \href{http://nimahejazi.org}
      {\tt \scriptsize \color{foreground} nimahejazi.org}
  \\[-4pt]
  \href{https://twitter.com/nshejazi}
      {\tt \scriptsize \color{foreground} twitter/@nshejazi}
  \\[-4pt]
  \href{https://github.com/nhejazi}
      {\tt \scriptsize \color{foreground} github/nhejazi}
}

%%%%%%%%%%%%%%%%%%%%%%%%%%%%%%%%%%%%%%%%%%%%%%%%%%%%%%%%%%%%%%%%%%%%%%%%%%%%%%%%

\begin{document}

% title slide
{
\setbeamertemplate{footline}{} % no page number here
\frame{
  \titlepage

  \vspace{-1em}

  \centerline{\href{https://goo.gl/8RWEy5}{\tt \scriptsize
                                           \underline{slides}: goo.gl/8RWEy5}}
  \vspace{-1.5em}
  \vfill \hfill \includegraphics[height=6mm]{Figs/cc-zero.png} \vspace*{-0.5cm}

  \note{This slide deck is for a reading group presentation on the manuscript
    ``Fair Inference on Outcomes'' (Rabi \& Shpitser, 2017), for the seminar on
    ``Fairness in Machine Learning'', organized in Fall 2017 by Moritz Hardt, at
    the University of California, Berkeley.

    Source: {\tt https://github.com/nhejazi/talk\_fair-outcomes} \\
    Slides: {\tt https://goo.gl/i3CxL9} \\
    With notes: {\tt https://goo.gl/8RWEy5}
}
}
}

%%%%%%%%%%%%%%%%%%%%%%%%%%%%%%%%%%%%%%%%%%%%%%%%%%%%%%%%%%%%%%%%%%%%%%%%%%%%%%%%

\begin{frame}[c]{Preview: Summary}
\only<1>{\addtocounter{framenumber}{-1}}

\begin{center}
\begin{itemize}
  \itemsep6pt
  \item Mediation analysis provides a framework under which intuitive
    definitions of fairness may be expressed.
  \item ``Fair inference'' is analogous to causal inference, except in that the
    counterfactuals explored refer to a ``fair'' world (n.~b., intentionally
    vague).
  \item Fairness may be characterized as the absence (or dampening) of a
    \textbf{path-specific effect (PSE)}.
  \item Restriction of a PSE is easily expressed as a likelihood maximization
    problem that features contraining the magnitude of the undesirable PSE.
  \item This approach to fairness avoids throwing away information (i.e.,
    ``fairness through unawareness'') but leaves the definition of fairness to
    the analyst.
\end{itemize}
\end{center}

\note{
We'll go over this summary again at the end of the talk. Hopefully, it will make
more sense then.
}

\end{frame}

%%%%%%%%%%%%%%%%%%%%%%%%%%%%%%%%%%%%%%%%%%%%%%%%%%%%%%%%%%%%%%%%%%%%%%%%%%%%%%%%

\begin{frame}[c]{Preliminaries: Notation}

\begin{center}
\begin{itemize}
  \itemsep10pt
  \item Data $\D = (Y, \bm{X})$; outcome $Y$ and feature vector $\bm{X}$.
  \item Sensitive features: $S \in \bm{X}$, where inference on $Y$ using $S$
    \textit{might} result in discrimination.
  \item Treatment variable: $A \in \bm{X}$.
  \item Mediator variables: $M \in \bm{X}$ or $\bm{M} \subseteq \bm{X}$.
  \item Potential outcome: $Y(a)$, realization of $Y$ under $A = a$.
\end{itemize}
\end{center}

\note{
}

\end{frame}

%%%%%%%%%%%%%%%%%%%%%%%%%%%%%%%%%%%%%%%%%%%%%%%%%%%%%%%%%%%%%%%%%%%%%%%%%%%%%%%%

\begin{frame}[c]{Preliminaries: Mediation Analysis}

\begin{center}
\begin{itemize}
  \itemsep10pt
  \item \textbf{Goal:} understand the mechanism by which $A$ influences $Y$.
  \item Decompose the \textbf{ACE} into \textit{direct} and \textit{indirect}
    effects mediated by a variable $M$.
  \item Partition feature space $\bm{X}$ into $A$ (treatment), $M$ (mediator),
    and $\bf{C} = \bm{X}$ \textbackslash \hspace{0.05em} $\{A, M\}$ (baseline
    factors).
  \item Counterfactual contrasts are expressed via \textit{nested} potential
    outcomes (i.e., $Y(a, M(a'))$).
\end{itemize}
\end{center}

\note{
  \begin{itemize}
    \item Nested potential outcomes read as ``the outcome $Y$ if $A$ were set
      to $a$ while $M$ were set to whatever value it would have attained had
      $A$ been set to $a'$''.
  \end{itemize}
}

\end{frame}

%%%%%%%%%%%%%%%%%%%%%%%%%%%%%%%%%%%%%%%%%%%%%%%%%%%%%%%%%%%%%%%%%%%%%%%%%%%%%%%%

\begin{frame}[c]{The Average Causal Effect (ACE)}

\begin{center}
\begin{itemize}
  \itemsep10pt
  \item $\text{ACE} = \E[Y(a)] - \E[Y(a')]$
  \item Not computed via $\E[Y \mid A],$ as associations between $A$ and $Y$ may
    be ``partly causal'' or spurious.
  \item Decomposition: $\text{ACE} = \text{NDE} + \text{NIE},$ where
    \textbf{NDE} is the \textit{Natural Direct Effect} and \textbf{NIE} is the
    \textit{Natural Indirect Effect}.
\end{itemize}
\begin{equation*}
\begin{split}
  \text{ACE} & = \E[Y(a)] - \E[Y(a')] \\ & = \E[Y(a)] - \E[Y(a, M(a')] \\ & +
    \E[(Y(a, M(a')] - \E[Y(a')]
\end{split}
\end{equation*}
\end{center}

\note{
  \begin{itemize}
    \item Decomposition of the ACE gives us a way to express undesirable PSEs
      using mediators
    \item The decomposition is just by way of a telescoping sum arguemnt.
  \end{itemize}
}

\end{frame}

%%%%%%%%%%%%%%%%%%%%%%%%%%%%%%%%%%%%%%%%%%%%%%%%%%%%%%%%%%%%%%%%%%%%%%%%%%%%%%%%

\begin{frame}[c]{The Natural \textit{Direct} Effect (NDE)}

\begin{center}

\begin{itemize}
  \itemsep10pt
  \item Comparison of the mean outcome under only the part of the treatment that
    directly affects it ($A = a$) and the placebo treatment (i.e., $A = a'$).
  \item Note that the \textit{indirect} effect of the treatment (through the
    mediator $M$) is ``turned off'' (i.e., $M(A = a')$).
\end{itemize}

\vspace{1em}
\begin{definition}{Natural \textbf{Direct} Effect}
$$\text{NDE} = \E[(Y(a, M(a')] - \E[Y(a')]$$
\end{definition}

\end{center}

\note{
With additional causal assumptions, the NDE is identified as the

\begin{definition}{Mediation formula}
$$
\sum_{\bm{C}, M} (\E[Y \mid a, M, \bm{C}] - \E[Y \mid a', M, \bm{C}])
  p(M \mid a', \bm{C}) p(\bm{C})
$$
\end{definition}

\begin{itemize}
    \item Estimation may be performed using plug-in estimators.
\end{itemize}
}

\end{frame}

%%%%%%%%%%%%%%%%%%%%%%%%%%%%%%%%%%%%%%%%%%%%%%%%%%%%%%%%%%%%%%%%%%%%%%%%%%%%%%%%

\begin{frame}[c]{The Natural \textit{Indirect} Effect (NIE)}

\begin{center}

\begin{itemize}
  \itemsep10pt
  \item Comparison of the outcome affected by all treatment (both direct and
    indirect) and the outcome where the effect through the mediator ($M$) is
    ``turned off'' (i.e., $M(A = a')$).
  \item Although in a roundabout manner, this quantity gets at the effect of the
    path-specific effect through the mediator on the outcome.
\end{itemize}

\vspace{1em}
\begin{definition}{Natural \textbf{Indirect} Effect}
$$\text{NIE} = \E[Y(a)] - \E[(Y(a, M(a')]$$
\end{definition}

\end{center}

\note{
}

\end{frame}

%%%%%%%%%%%%%%%%%%%%%%%%%%%%%%%%%%%%%%%%%%%%%%%%%%%%%%%%%%%%%%%%%%%%%%%%%%%%%%%%

\begin{frame}[c]{Example: \textit{Thank You for Smoking}}

For a better intuition of $Y(a, M(a'))$, consider the following:

\begin{center}
\begin{itemize}
  \itemsep4pt
  \item Let $Y$ be a health outcome (e.g., survival probability), $A$ be a
    treatment (e.g., smoking).
  \item Consider a decomposition of the effect of $A$ on $Y$ --- that is, let
    $M$ be a mediator (e.g., cancer).
  \item $A$ affects $Y$ directly (nicotine exposure) and indirectly (inducing
    lung cancer, through $M$).
  \item Here, $Y(a, M(a'))$ correponds to ``the response of $Y$ to an
    intervention that sets the nicotine exposure (direct effect) to what it
    would be in smokers, and the smoke exposure (indirect effect) to what it
    would be in non-smokers'' (e.g., nicotine patch).
\end{itemize}
\end{center}

\note{
}

\end{frame}

%%%%%%%%%%%%%%%%%%%%%%%%%%%%%%%%%%%%%%%%%%%%%%%%%%%%%%%%%%%%%%%%%%%%%%%%%%%%%%%%

\begin{frame}[c]{Path-Specific Effects}

\begin{center}
\begin{itemize}
  \itemsep10pt
  \item A more general idea than the NDE and NIE --- such effects are easily
    formulated as nested counterfactuals.
  \item \textit{Intuition}: along a path of interest, all nodes behave as if the
    active rule were imposed (i.e., $A = a$) while, along all other paths, nodes
    behave as though the alternative were the case (i.e., $A = a'$).
\end{itemize}

\vspace{1em}
\begin{definition}{Path-Specific Effect (PSE)}
\\
\vspace{1em}
(Along a path, say $A \rightarrow W \rightarrow Y$)
$$\E[Y(a', W(M(a'), a), M(a'))] - \E[Y(a')]$$
\end{definition}

\end{center}

\note{
With more complex assumptions, we get the
\begin{definition}{Edge G-formula}
$$
\sum_{\bm{C}, M, W} \E[Y \mid a', W, M, \bm{C}] p(W \mid a, M, \bm{C})
  p(M \mid a', \bm{C}) p(\bm{C})
$$
\end{definition}

\begin{itemize}
    \item Estimation may be performed using plug-in estimators.
\end{itemize}
}

\end{frame}

%%%%%%%%%%%%%%%%%%%%%%%%%%%%%%%%%%%%%%%%%%%%%%%%%%%%%%%%%%%%%%%%%%%%%%%%%%%%%%%%

\begin{frame}[c]{Finding Fairness}

\begin{center}
\begin{itemize}
  \itemsep4pt
  \item Much work has focused on defining fairness via associative relationships
    (including equalized odds). Such criteria provided unintuive results when
    the sensitive feature is not randomly assigned.
  \item Here, an approach that ought to provide intuitive results (wrt
    fairness), even when the sensitive attribute is associated with the outcome
    (perhaps by way of an unobserved feature), is proposed.
  \item Associative fairness metrics fail to properly model sources of
    confounding (between $S$ and $Y$).
  \item Generally, this failure is rooted in the fact that ``counterfactual
    probabilities are complex functions of the observed data, no just
    conditional densities.''
\end{itemize}
\end{center}

\note{
To see the problem with associative fairness criteria, consider the following
example (see section 4, paragraph 3 of paper for details):
\begin{itemize}
  \itemsep10pt
  \item Letting $p(H \mid C, G)$ be a hiring rule, we would assess fairness as
    follows: $p(H = 1 \mid C = 1) = p(H = 1 \mid C = 0) \approx 0.022$, using
    associative fairness criteria.
  \item Intuitively, fairness of a hiring rule would lead to equal hiring
    probabilities for cases and controls in a hypothetical randomized trial
    (RCT) --- i.e., $p(H(C = 1)) = p(H(C = 0))$.
  \item In our example, application of the adjustment formula would yield $p(H(C
    = 1)) = 0.025$ and $p(H(C = 0)) = 0.25$, a rather striking difference.
  \item The difference between $p(H(C = 0))$ and $p(H \mid C = 0)$ is driven by
    extreme values of $p(C \mid G)$.
  \item This is basically Simpson's paradox, or rather close to it.
\end{itemize}
}

\end{frame}

%%%%%%%%%%%%%%%%%%%%%%%%%%%%%%%%%%%%%%%%%%%%%%%%%%%%%%%%%%%%%%%%%%%%%%%%%%%%%%%%

\begin{frame}[c]{In Pursuit of ``Fair Inference''}

\begin{center}
\begin{itemize}
  \itemsep10pt
  \item Fairness is, at its core, rooted in counterfactuals. Thus, we can see
    \textit{``fair inference''} as a branch of causal inference wherein the
    counterfactuals to be considered are with respect to a ``fair'' world.
  \item \textit{Discrimination} may be expressed as the presence of a particular
    PSE, with choice of the specific PSE left as a domain-specific issue.
  \item Thus, minimization of specific PSEs corresponds to minimizing
    discrimination and is a problem of constrained inference on statistical
    models.
\end{itemize}
\end{center}

\note{
\textit{From the legal literature:} \\
``The central question in any employment-discrimination case is whether the
employer would have taken the same action had the employee been of a different
race (age, sex, religion, national origin, etc.) and everything else had been
the same.''
}

\end{frame}

%%%%%%%%%%%%%%%%%%%%%%%%%%%%%%%%%%%%%%%%%%%%%%%%%%%%%%%%%%%%%%%%%%%%%%%%%%%%%%%%

\begin{frame}[c]{Fairness as PSE Minimization}

\begin{center}
\begin{itemize}
  \itemsep10pt
  \item Let $p(Y, \bf{X})$ be a statistical model, assumed to be induced by a
    \textit{causal model}.
  \item Discrimination (wrt $Y$ based on $S \in \bm{X}$) in this model is a PSE,
    identified as the functional $f(p(Y, \bm{X}))$.
  \item Let $(\epsilon_l, \epsilon_u)$ be lower and upper bounds on the PSE,
    giving the degree of unfairness considered tolerable (n.b., the PSE is
    removed in the special case $\epsilon_l = \epsilon_u$).
  \item \underline{\textbf{Proposal}}: transform $p(Y, \bf{X})$ into $p^*(Y,
    \bm{X})$ under the constraint that the PSE of interest lies within
    $(\epsilon_l, \epsilon_u)$, where the two distributions are close in the
    sense of KL-divergence.
\end{itemize}
\end{center}

\note{
\begin{definition}{Kullback-Leibler Divergence}
$$\infdiv{P}{Q} = \int_{-\infty}^{\infty} p(x) \text{log} \frac{p(x)}{q(x)}
  \text{dx}$$
\end{definition}

\begin{itemize}
  \item Why do we restrict ourselves to KL-divergence?
  \item Well, as we'll see shortly, our approach amounts to a constrained
    likelihood maximization problem. This makes KL-divergence a natural choice.
  \item Were we to consider formulating our fairness metric differently, we
    could use other distance metrics (e.g., Hellinger distance, total variation
    distance).
\end{itemize}
}

\end{frame}

%%%%%%%%%%%%%%%%%%%%%%%%%%%%%%%%%%%%%%%%%%%%%%%%%%%%%%%%%%%%%%%%%%%%%%%%%%%%%%%%

\begin{frame}[c]{Finding Fair Worlds I}

\begin{center}
\begin{itemize}
  \itemsep10pt
  \item \underline{\textbf{Proposal}}: We can make \textit{any} function of $p$
    \textbf{\textit{fair}}, merely by computing it from $p^*$ (instead of from
    $p$).
  \item To ensure fairness, we must make inference only in the ``fair world'',
    just as we only perform inference on counterfactuals in causal inference.
  \item To do this, map any $x^i$ from $p$ to a sensible version of it drawn
    from $p^*$ --- i.e., find a $g: x^i_{p} \mapsto x^i_{p^*}$.
  \item I want to be fair, so what exactly do I do?
\end{itemize}
\end{center}

\note{
\begin{itemize}
  \item The intuition here is simply that $p^*$ includes all the same
    information that was found in $p$, except that information that leads to
    discrimination (is used in the unfavorable PSE).
  \item We'll shortly see what it means to construct a fair world $p^*$.
\end{itemize}
}

\end{frame}

%%%%%%%%%%%%%%%%%%%%%%%%%%%%%%%%%%%%%%%%%%%%%%%%%%%%%%%%%%%%%%%%%%%%%%%%%%%%%%%%

\begin{frame}[c]{Finding Fair Worlds II}

\begin{center}
\begin{itemize}
  \itemsep10pt
  \item Consider the following general setup:
    \begin{itemize}
      \itemsep3pt
      \item finite samples $\D$ drawn from $p(Y, \bm{X})$
      \item a likelihood function $\lik_{Y, \bm{X}}(\D; \alpha)$
      \item a discriminative PSE $f(p(Y, \bf{X}))$ with bounds $(\epsilon_l,
        \epsilon_u)$
      \item an estimator of the PSE $g(\D)$.
    \end{itemize}
  \item We obtain fairness by solving:
    $$\hat{\alpha} = \argmax_{\alpha} \lik_{Y, \bm{X}} (\D; \alpha),$$
    subject to $\epsilon_l \leq g(\D) \leq \epsilon_u.$
  \item In this setup, fairness is achieved by constraining parts of $p(Y,
    \bm{X},; \alpha),$ with the choice of $g$ determining exactly what is
    constrained.
\end{itemize}
\end{center}

\note{
\begin{itemize}
  \item Note that knowing $p$ and $p^*$ exactly would make the problem rather
    easy, since we would merely use $x_{\bm{W}}^i,$ where $\bm{W}$ is simply the
    largest subset of $X$ that contrains the PSE in a desirable manner (i.e.,
    where $p(\bm{W}) = p^*(\bm{W})$.
  \item Further, it's important to consider ``large'' statistical models ---
    perhaps infinite-dimensional models? --- in order to ensure that prediction
    (and inference, if so desired) can be optimized with respect to
    out-of-sample observations.
  \item In regard to both prediction and inference, it turns out that most
    estimators that rely on the outcome ($Y$) model are not robust to its
    misspecification. Thus, it is important to either use simpler (IPW)
    estimators or more modern approaches (e.g., \textit{triply robust}
    estimators).
\end{itemize}
}

\end{frame}

%%%%%%%%%%%%%%%%%%%%%%%%%%%%%%%%%%%%%%%%%%%%%%%%%%%%%%%%%%%%%%%%%%%%%%%%%%%%%%%%

\begin{frame}[c]{Fairness is (Partial?) (Un)Awareness}

\begin{center}
\begin{itemize}
  \itemsep10pt
  \item Since using all of the information contained in $p$ leads to unfairness,
    this approach amounts to discarding information that is exclusively in $p$,
    relative to $p^*$.
  \item The goal of this approach is to use the available information as well as
    possible, but only in so far as our inferences are drawn from the ``fair
    world.''
  \item In this approach, fairness is characterized as the \textit{a priori}
    \underline{inadmissability} of certain paths in the DAG of interest --- that
    is, paths other than a single edge path might cause discrimination.
\end{itemize}
\end{center}

\note{
\begin{itemize}
  \item Consider the experiment of the authors using the ``adult'' dataset:
    here, there is concern about the direct effect of gender on income class as
    well as the effect of gender on income class through marital status. Using
    the proposed constrained optimization approach, they are able to avoid
    dropping sensitive variables from the analysis while improving upon the
    efficacy of ``blind'' models.
  \item Specifically, they estimate the PSE in the unconstrained model to be
    $3.16$, and then restric the PSE to lie within $(0.95, 1.00)$ in the
    ``fair'' (constrained) model. The ``fair'' model achives accuracy of $72\%$
    while the ``blind'' model only achieves $42\%$ (the unconstrained model does
    better, with an accuracy of $82\%$).
\end{itemize}
}

\end{frame}

%%%%%%%%%%%%%%%%%%%%%%%%%%%%%%%%%%%%%%%%%%%%%%%%%%%%%%%%%%%%%%%%%%%%%%%%%%%%%%%%

\begin{frame}[c]{Review: Summary}

\begin{center}
\begin{itemize}
  \itemsep6pt
  \item Mediation analysis provides a framework under which intuitive
    definitions of fairness may be expressed.
  \item ``Fair inference'' is analogous to causal inference, except in that the
    counterfactuals explored refer to a ``fair'' world (n.~b., intentionally
    vague).
  \item Fairness may be characterized as the absence (or dampening) of a
    \textbf{path-specific effect (PSE)}.
  \item Restriction of a PSE is easily expressed as a likelihood maximization
    problem that features contraining the magnitude of the undesirable PSE.
  \item This approach to fairness avoids throwing away information (i.e.,
    ``fairness through unawareness'') but leaves the definition of fairness to
    the analyst.
\end{itemize}
\end{center}

\note{
It's always good to include a summary. You've seen this all before.
}

\end{frame}

%%%%%%%%%%%%%%%%%%%%%%%%%%%%%%%%%%%%%%%%%%%%%%%%%%%%%%%%%%%%%%%%%%%%%%%%%%%%%%%%

%% don't want dimming of references
\setbeamercovered{}
\beamerdefaultoverlayspecification{}

\begin{frame}[c,allowframebreaks]{References}

\bibliographystyle{apalike}
\nocite{*}
\bibliography{references}

%\note{Here's some work we've talked about. Go check these out if interested.}

\end{frame}

%%%%%%%%%%%%%%%%%%%%%%%%%%%%%%%%%%%%%%%%%%%%%%%%%%%%%%%%%%%%%%%%%%%%%%%%%%%%%%%%

\begin{frame}[c]{Thank you.}

\Large
Slides: \href{https://goo.gl/i3CxL9}{goo.gl/i3CxL9} \quad
\includegraphics[height=5mm]{Figs/cc-zero.png}

\vspace{3mm}
Notes: \href{https://goo.gl/8RWEy5}{goo.gl/8RWEy5}

\vspace{3mm}
Source (repo): \href{https://goo.gl/qJSoz6}{goo.gl/qJSoz6}

\vspace{3mm}
\href{https://www.stat.berkeley.edu/~nhejazi}{\tt
  stat.berkeley.edu/\textasciitilde{}nhejazi}

\vspace{3mm}
\href{http://nimahejazi.org}{\tt nimahejazi.org}

\vspace{3mm}
\href{https://twitter.com/nshejazi}{\tt twitter/@nshejazi}

\vspace{3mm}
\href{https://github.com/nhejazi}{\tt github/nhejazi}

\note{Here's where you can find me, as well as the materials from this talk.}

\end{frame}

%%%%%%%%%%%%%%%%%%%%%%%%%%%%%%%%%%%%%%%%%%%%%%%%%%%%%%%%%%%%%%%%%%%%%%%%%%%%%%%%

\end{document}

